\documentclass[11pt,a4paper]{article}

\usepackage[margin=2.5cm]{geometry}
\usepackage{amsmath}
\usepackage{amssymb}
\usepackage{booktabs}
\usepackage{siunitx}
\usepackage{hyperref}

\title{Time Synchronisation Between Klipper Motion and Trinkey Accelerometer Data}
\author{}
\date{}

\begin{document}

\maketitle

\section{Purpose}

The frequency response function (FRF) identification requires the commanded
input acceleration and the measured output acceleration to be expressed on the
same time axis. In this experiment, the input motion is generated by Klipper,
while the output acceleration is measured by two BNO055 accelerometers connected
to a separate Trinkey QT2040 board. Therefore, the system contains two relevant
time bases:

\begin{itemize}
    \item Klipper \emph{print time}, which is the time base used to schedule
    motion on the printer MCU.
    \item Trinkey time, given by the timestamp \(t_{\mathrm{us}}\) produced by
    the RP2040 timer on the Trinkey.
\end{itemize}

The aim of the synchronisation procedure is to map each Trinkey timestamp to
Klipper print time. After this mapping, the input and output signals can be
compared sample-by-sample using a common time coordinate.

\section{Klipper Print Time}

Klipper plans motion on the host computer, but the final step pulses are
executed by the printer MCU. Internally, Klipper uses a variable called
\emph{print time} to represent the scheduled time of a motion command. For the
primary MCU, print time is directly related to the MCU clock:

\begin{equation}
    c = f_{\mathrm{mcu}} p,
\end{equation}

where \(c\) is the MCU clock value in ticks, \(f_{\mathrm{mcu}}\) is the MCU
clock frequency, and \(p\) is Klipper print time in seconds. Therefore, print
time can be interpreted as the MCU clock expressed in seconds. Klipper's code
overview documents this host/MCU clock model and the conversion between MCU
clock values and print time~\cite{klippercodeoverview}.

When a vibration test is executed, each input motion segment has a known print
time. If the first segment starts at print time \(p_0\), and a generated segment
starts at local test time \(\tau_k\), then its scheduled start time is

\begin{equation}
    p_k = p_0 + \tau_k.
\end{equation}

The corresponding input acceleration for that segment is known from the
trajectory generation code. Therefore, the input acceleration can be stored as a
piecewise constant signal in Klipper print time:

\begin{equation}
    u(p) = a_k, \qquad p_k \leq p < p_k + \Delta t_k,
\end{equation}

where \(a_k\) is the acceleration of segment \(k\), and \(\Delta t_k\) is its
duration.

\section{Trinkey Time}

The Trinkey firmware records accelerometer samples with its own timestamp
\(t_{\mathrm{us}}\), measured in microseconds. This timestamp is produced by the
Trinkey RP2040 and is independent from the printer MCU clock. For analysis it is
first converted to seconds:

\begin{equation}
    q = 10^{-6} t_{\mathrm{us}}.
\end{equation}

The two BNO055 accelerometers are connected to the same Trinkey board. Therefore,
both accelerometers share the same Trinkey time base. This is useful, because
one time synchronisation model can be used for both the toolhead and bed
accelerometer data.

\section{Clock Model}

Over the duration of one vibration sweep, the relation between Trinkey time and
Klipper print time can be modelled as an affine function:

\begin{equation}
    p = \alpha q + \beta,
    \label{eq:clock_model}
\end{equation}

where:

\begin{itemize}
    \item \(q\) is Trinkey time in seconds,
    \item \(p\) is Klipper print time in seconds,
    \item \(\alpha\) is the relative clock rate,
    \item \(\beta\) is the clock offset.
\end{itemize}

If both clocks ran at exactly the same rate, then \(\alpha = 1\). In practice,
\(\alpha\) can be slightly different from one due to oscillator tolerances and
clock drift. The offset \(\beta\) accounts for the fact that the two clocks do
not start from the same zero time. This is the standard clock offset and clock
skew model used in clock synchronisation methods, including sensor-network
synchronisation approaches~\cite{elson2002rbs}.

\section{Synchronisation Measurements}

To estimate \(\alpha\) and \(\beta\), the Klipper host sends repeated
synchronisation requests to the Trinkey. The vibration-test logger takes 20
observations before and 20 after each excitation. Continuous streaming instead
maintains a rolling window of at most 64 observations and normally requests one
new observation per second. For one synchronisation request:

\begin{enumerate}
    \item The Klipper host records its monotonic time before sending the request:
    \(h_{0,i}\).
    \item The host sends a synchronisation command to the Trinkey.
    \item The Trinkey immediately replies with its current timestamp
    \(q_i\).
    \item The host records its monotonic time after receiving the reply:
    \(h_{1,i}\).
\end{enumerate}

The host time corresponding to the Trinkey reply is estimated by the midpoint:

\begin{equation}
    h_i = \frac{h_{0,i} + h_{1,i}}{2}.
    \label{eq:midpoint}
\end{equation}

The round-trip time is

\begin{equation}
    r_i = h_{1,i} - h_{0,i}.
\end{equation}

The midpoint approximation assumes that the outward and return communication
delays are approximately symmetric. This is not perfectly true for USB serial
communication, but repeated measurements allow observations with additional
queueing delay to be rejected. This follows the general principle of two-way
clock synchronisation methods such as Cristian's algorithm and the Network Time
Protocol (NTP), where round-trip delay is used to estimate and qualify a clock
offset observation~\cite{cristian1989probabilistic,mills2010ntp}.

The host monotonic time \(h_i\) is then converted to Klipper print time using
Klipper's clock synchronisation estimate:

\begin{equation}
    p_i = \operatorname{estimated\_print\_time}(h_i).
\end{equation}

Each synchronisation request therefore gives one calibration pair:

\begin{equation}
    (q_i, p_i).
\end{equation}

\section{Least-Squares Clock Fit}

The parameters \(\alpha\) and \(\beta\) are estimated by fitting the model in
Equation~\eqref{eq:clock_model} to all accepted synchronisation pairs. This is
an ordinary linear least-squares (OLS) problem with an intercept
term~\cite{montgomery2012linear}. The measurements can be written in matrix form as

\begin{equation}
    p = X\theta + e,
\end{equation}

where

\begin{equation}
    p =
    \begin{bmatrix}
        p_1 \\
        p_2 \\
        \vdots \\
        p_N
    \end{bmatrix},
    \qquad
    X =
    \begin{bmatrix}
        1 & q_1 \\
        1 & q_2 \\
        \vdots & \vdots \\
        1 & q_N
    \end{bmatrix},
    \qquad
    \theta =
    \begin{bmatrix}
        \beta \\
        \alpha
    \end{bmatrix}.
\end{equation}

Here, \(p\) is the vector of Klipper print times, \(q_i\) are the Trinkey
timestamps converted to seconds, and \(e\) is the timing error in the
synchronisation measurements. The first column of \(X\) is equal to one because
the model includes the offset \(\beta\). The least-squares estimate is obtained
by minimizing the squared residual norm:

\begin{equation}
    \hat{\theta}
    =
    \arg\min_{\theta} \left\|p - X\theta\right\|_2^2.
\end{equation}

The closed-form normal-equation solution is

\begin{equation}
    \hat{\theta} = (X^T X)^{-1} X^T p,
    \label{eq:matrix_ls}
\end{equation}

provided that \(X^T X\) is non-singular. In implementation, the inverse should
not be formed explicitly. The implementation uses MATLAB's backslash operator,
\(\hat{\theta} = X \backslash p\), which returns a least-squares solution for an
overdetermined rectangular system~\cite{mathworksmldivide}. This avoids explicitly
forming the normal-equation inverse.

After these parameters have been estimated, every accelerometer sample can be
mapped to Klipper print time:

\begin{equation}
    p_{\mathrm{sample}} =
    \alpha \left(10^{-6} t_{\mathrm{us}}\right) + \beta.
    \label{eq:sample_mapping}
\end{equation}

\section{Latency Rejection}

USB and operating system scheduling can occasionally delay a synchronisation
request. Such delayed requests should not be used directly, because they reduce
the accuracy of the clock fit. The round-trip time \(r_i\) is used as a quality
measure.

The selected rejection rule is

\begin{equation}
    r_i \leq r_{\min} + r_{\mathrm{tol}},
\end{equation}

where \(r_{\min}\) is the minimum round-trip time in the current set of
observations and the tolerance is fixed to

\begin{equation}
    r_{\mathrm{tol}} = \SI{1}{ms}.
\end{equation}

The OLS fit uses only the accepted index set

\begin{equation}
    \mathcal{A} = \{i : r_i \leq r_{\min} + \SI{1}{ms}\}.
\end{equation}

Giving preference to low-delay observations is consistent with established
clock-filter practice: for example, the NTPv4 clock filter orders observations
by increasing round-trip delay before selecting the synchronization
sample~\cite[Sec.~10]{mills2010ntp}. However, NTP does not prescribe the
\(\SI{1}{ms}\) margin used here. That value is an experiment-specific engineering
threshold for this USB link and must therefore be reported as part of the method,
along with the number of observations retained. It should be revalidated if the
host, USB connection, firmware, or sampling workload changes.

\section{Implemented Estimator}

The same rejection rule and unweighted OLS model are used by the online and
offline implementations:

\begin{itemize}
    \item The live mapper is implemented in
    \path{external/klipper/klippy/extras/trinkey_accel.py}. It keeps at most 64
    observations. During start-up it estimates only the offset. It enables the
    affine OLS fit after at least eight accepted observations span at least
    \SI{5}{s}. A fitted clock rate outside the sanity interval
    \(0.995 \leq \alpha \leq 1.005\) is rejected in favour of the offset-only
    estimate.
    \item Vibration-test analysis is implemented in
    \path{analysis/matlab/frf_align_signals.m}. It verifies that all 40 expected
    observations are present, applies the same \(r_{\min}+\SI{1}{ms}\) filter,
    requires at least three accepted observations including at least one before
    and one after the excitation, and then performs the affine OLS fit.
\end{itemize}

The online start-up guard is not a different clock model. It prevents a poorly
conditioned slope estimate when the accepted timestamps cover too little elapsed
time. The offline data already span the complete experiment, so that guard is not
needed there. The retained count, minimum RTT, RTT limit, tolerance, slope, and
offset are exposed or stored as diagnostics so the fitted mapping can be audited.

\section{Alignment of Input and Output Signals}

After synchronisation, both the input and output are represented in Klipper
print time. The input is the commanded acceleration \(u(p)\), obtained from the
motion segments generated by Klipper. The output is the measured acceleration
\(y(p)\), obtained by converting each Trinkey sample timestamp using
Equation~\eqref{eq:sample_mapping}.

The two signals can then be compared on the same time axis:

\begin{equation}
    u(p) \quad \text{and} \quad y(p).
\end{equation}

For frequency-domain identification, the steady-state part of the response is
selected after the initial ramp. The complex gain at angular frequency
\(\omega = 2\pi f\) can be estimated by fitting sine and cosine components to
the measured output:

\begin{equation}
    y(p) \approx B\sin(\omega p) + C\cos(\omega p).
\end{equation}

The same time base must be used for the input and output. Otherwise, a timing
error \(\Delta t\) introduces a phase error

\begin{equation}
    \Delta \phi = -\omega \Delta t.
\end{equation}

In degrees this is

\begin{equation}
    \Delta \phi_{\mathrm{deg}} =
    -360 f \Delta t.
\end{equation}

This shows why time synchronisation is important. For example, a timing error of
\(\SI{1}{ms}\) causes a phase error of \(\SI{1.8}{deg}\) at
\(\SI{5}{Hz}\), but \(\SI{36}{deg}\) at \(\SI{100}{Hz}\).

\section{Remaining Sources of Phase Error}

The proposed synchronisation method removes the main timing uncertainty caused
by USB serial latency and host-side logging. However, it does not remove all
possible phase errors. The following effects may still remain:

\begin{itemize}
    \item BNO055 internal filtering and sensor group delay.
    \item Small timing differences between reading the bed and toolhead
    accelerometers, because they are read sequentially on the same I2C bus.
    \item Quantisation and jitter in the generated motion segments.
    \item Mechanical phase lag caused by the printer structure, which is the
    desired quantity for FRF identification.
\end{itemize}

These effects should be considered when interpreting the measured phase. In
particular, any approximately constant additional delay can be estimated and
reported separately from the mechanical dynamics.

\section{Summary}

The synchronisation procedure maps Trinkey accelerometer timestamps onto
Klipper print time using repeated sync requests and a linear clock model. The
main steps are:

\begin{enumerate}
    \item Log the commanded input acceleration using Klipper print time.
    \item Log each accelerometer sample with its Trinkey timestamp.
    \item Collect synchronisation pairs between Trinkey time and Klipper print
    time.
    \item Retain pairs satisfying
    \(r_i \leq r_{\min}+\SI{1}{ms}\).
    \item Fit the affine model \(p = \alpha q + \beta\) by unweighted OLS using
    only the retained pairs.
    \item Convert all accelerometer samples to Klipper print time.
    \item Use the common time base for FRF estimation.
\end{enumerate}

This approach follows the same principle as Klipper's built-in accelerometer
handling: sample timestamps are not based on when data arrives at the host, but
on a local sensor clock that is later translated to Klipper's motion time
base~\cite{klipperapiserver}.

\begin{thebibliography}{9}

\bibitem{cristian1989probabilistic}
F.~Cristian,
\newblock ``Probabilistic clock synchronization,''
\newblock \emph{Distributed Computing}, vol.~3, no.~3, pp.~146--158, 1989.
\newblock doi: \href{https://doi.org/10.1007/BF01784024}{10.1007/BF01784024}.

\bibitem{mills2010ntp}
D.~L. Mills, J.~Martin, J.~Burbank, and W.~Kasch,
\newblock ``Network Time Protocol Version 4: Protocol and Algorithms Specification,''
\newblock RFC 5905, Internet Engineering Task Force, 2010.
\newblock doi: \href{https://doi.org/10.17487/RFC5905}{10.17487/RFC5905}.

\bibitem{elson2002rbs}
J.~Elson, L.~Girod, and D.~Estrin,
\newblock ``Fine-grained network time synchronization using reference broadcasts,''
\newblock in \emph{Proceedings of the 5th Symposium on Operating Systems Design
and Implementation}, pp.~147--163, 2002.
\newblock doi: \href{https://doi.org/10.1145/844128.844143}{10.1145/844128.844143}.

\bibitem{montgomery2012linear}
D.~C. Montgomery, E.~A. Peck, and G.~G. Vining,
\newblock \emph{Introduction to Linear Regression Analysis},
\newblock 5th~ed. Wiley, 2012.
\newblock ISBN: \href{https://catalogimages.wiley.com/images/db/pdf/9780470542811.excerpt.pdf}{978-0-470-54281-1}.

\bibitem{mathworksmldivide}
MathWorks,
\newblock ``\texttt{mldivide}: Solve systems of linear equations
\(A x = B\) for \(x\),'' MATLAB documentation.
\newblock \url{https://www.mathworks.com/help/matlab/ref/double.mldivide.html}.

\bibitem{klippercodeoverview}
Klipper Project,
\newblock ``Code overview: Time,'' Klipper documentation.
\newblock \url{https://www.klipper3d.org/Code_Overview.html}.

\bibitem{klipperapiserver}
Klipper Project,
\newblock ``API server: \texttt{adxl345/dump\_adxl345},'' Klipper documentation.
\newblock \url{https://www.klipper3d.org/API_Server.html}.

\end{thebibliography}

\end{document}
